\documentclass[10.5pt,a4paper]{article}

\usepackage[top=0.85in,bottom=0.85in,left=0.9in,right=0.9in]{geometry}
\usepackage[T1]{fontenc}
\usepackage{lmodern}
\usepackage{microtype}
\usepackage{parskip}
\usepackage{booktabs}
\usepackage{array}
\usepackage{tabularx}
\usepackage{enumitem}
\usepackage[hidelinks]{hyperref}

\title{\textbf{BioGuard: A Portable Conversation-Layer Biosecurity Protocol}\\[0.2em]\normalsize AIxBio Hackathon Submission Draft}
\author{Jason Tang}
\date{}

\begin{document}
\maketitle
\vspace{-1em}

\section*{Abstract}

Biosecurity controls still concentrate on downstream sequence screening or model-level policies, but multi-turn AI conversations can transfer dangerous biological design knowledge before either layer is triggered. BioGuard is a hackathon-scoped proof of a protocol ecosystem seed. It packages a portable conversation-layer control path as reusable \texttt{SKILL.md} blocks, backed by a Biological Knowledge Transfer (BKT) scoring contract and a deterministic evaluation path. The hackathon package aims for an ambitious but disciplined proof: one deep host integration, additional host compatibility demonstrations, a stronger adversarial benchmark slice, baseline and ablation comparisons, and a reproducibility trail across contract versions and evidence artifacts. The output is therefore a protocol object with explicit conformance and distribution evidence, not a mere prototype.

\section*{Problem}

The dangerous transaction is the conversation. Existing biosecurity infrastructure still assumes risk is best measured at synthesis order time or by downstream model-output checks. Under stronger AI capability conditions, that assumption becomes less defensible. Multi-turn dialogue can gradually assemble a usable biological design protocol without presenting an obviously hazardous single prompt, sequence, or order request. The missing control layer is therefore the bio-LLM interface itself.

\section*{Hackathon Claim}

The hackathon claim is intentionally disciplined but no longer minimal: \textbf{a portable conversation-layer protocol ecosystem, backed by the BKT scoring contract and a deployable Skills Profile, can detect dangerous biological knowledge transfer on multi-turn cases better than stateless controls while remaining portable across multiple host surfaces}. This is a protocol claim with evidence, not a platform launch claim.

\section*{System}

BioGuard has three core runtime skills and one optional extension:
\begin{itemize}[leftmargin=1.2em]
  \item \textbf{bioguard-bkt-scoring:} canonical event-level BKT scoring logic.
  \item \textbf{bioguard-bio-trace:} context-aware tracing over a short conversation window.
  \item \textbf{bioguard-bio-guard:} operator-facing wrapper that returns a final decision envelope and audit trace.
  \item \textbf{bioguard-bio-seq (optional):} sequence-risk extension used only for ablation or future work.
\end{itemize}

These skills are packaged as \texttt{SKILL.md} blocks on top of the existing Agent Skills ecosystem \cite{agentskills2026,claude2026,githubclispec,copilot2026,openai2026}. BioGuard does not propose a new general skills standard; it defines a narrow profile and a shared scoring contract for a specific biosecurity use case.

The benchmark surface is separate. \textbf{BioSession} is the adversarial corpus used to test the protocol. It is not a runtime skill.

\section*{Conformance and Standards Core}

The protocol ecosystem centers on three required artifacts beyond model performance:
\begin{itemize}[leftmargin=1.2em]
  \item host capability profile (what is implemented, unsupported, and fallback),
  \item host deviation matrix (differences from protocol defaults),
  \item conformance report (human-readable portability decisions and risks).
\end{itemize}

A successful submission must show these artifacts as reviewable proof that portability is being measured, not only claimed.

\section*{Evaluation}

The primary comparison is against stateless baselines on multi-turn cases. The primary metric is recall at a fixed 5\% false positive rate, reported with bootstrap confidence intervals. The comparison set should include:
\begin{itemize}[leftmargin=1.2em]
  \item BioGuard composite screening,
  \item keyword-filter baseline,
  \item Llama Guard 3,
  \item GPT-5.4 zero-shot,
  \item pre-inference-only ablation,
  \item post-inference-only ablation,
  \item and host portability/deviation results.
\end{itemize}

Success, for hackathon purposes, means a clean and reproducible result with interpretable outputs, a stronger-than-baseline empirical story or deeply informative failure, and enough host evidence to make portability feel real rather than promised. Partial success still matters if BioGuard produces a clear taxonomy of misses and a concrete protocol revision path.

\section*{Scope}

The ambitious hackathon package is constrained to what can be defended rigorously while still aiming at a top-tier result:
\begin{itemize}[leftmargin=1.2em]
  \item one locked BKT contract,
  \item one deep host-compatible skill path,
  \item two additional host compatibility demonstrations,
  \item one stronger deterministic benchmark slice,
  \item one primary comparison table and one ablation table,
  \item one portability matrix and operator replay path,
  \item and one reproducibility trail plus dual-use appendix.
\end{itemize}

Out of scope for the hackathon package are treating BioSeq as the main contribution, making open-weight fine-tuning a success condition, or drifting into a general CBRN platform thesis.

\section*{Why This Is Compelling}

The strongest attractor is architectural: BioGuard moves the safety boundary to the layer where dangerous knowledge transfer actually occurs. The second is practical: the project uses a real, already emerging Agent Skills substrate rather than a hypothetical future interoperability story. The third is operational: the protocol yields auditable, replayable decisions rather than opaque scores. The fourth is evaluative: the package is designed to produce a legible result even if performance uplift is partial or null.

\section*{Limitations and Dual-Use Considerations}

This package still has clear limitations. Session-scale behavior may outrun turn-window scoring, host compatibility does not by itself prove full semantic portability, and stronger benchmark coverage does not eliminate false positives or false negatives. Because the work touches biosecurity screening, public artifacts should avoid disclosing operational misuse details or raw hazardous biological content beyond what is necessary to explain the evaluation. If portability or policy weaknesses are discovered in a concrete host environment, disclosure should prioritize maintainers and operators before public amplification. The ethical bar for this project is decision-useful evidence with explicit uncertainty, not inflated deployment claims.

\section*{Governance Commitments}

For this track, governance includes:
\begin{itemize}[leftmargin=1.2em]
  \item explicit versioned contracts (\texttt{spec/*}),
  \item explicit handling of unsupported-host features,
  \item reproducible evidence commands,
  \item and protocol revision notes when conformance findings require scope changes.
\end{itemize}

\section*{Next Step}

If the hackathon package succeeds, the fellowship-scale follow-up is straightforward: expand BioSession, harden rubric reliability, deepen host conformance, and publish a full empirical report on where the protocol helps, where it fails, and how the boundary should be revised.

\section*{References}

\begin{thebibliography}{99}

\bibitem{agentskills2026}
Agent Skills. 2026.
\textit{The SKILL.md Open Standard for Reusable Agent Capabilities.}
\url{https://agentskills.io}

\bibitem{claude2026}
Anthropic. 2026.
\textit{Claude Code Skills.}
\url{https://code.claude.com/docs/en/skills}

\bibitem{githubclispec}
GitHub. 2026.
\textit{Manage agent skills with GitHub CLI.}
\url{https://github.blog/changelog/2026-04-16-manage-agent-skills-with-github-cli/}

\bibitem{copilot2026}
GitHub. 2026.
\textit{Add skills to cloud agents.}
\url{https://docs.github.com/en/copilot/how-tos/copilot-on-github/customize-copilot/customize-cloud-agent/add-skills}

\bibitem{openai2026}
OpenAI. 2026.
\textit{Codex and Agent Skills.}
\url{https://openai.com/codex/}

\bibitem{ibbis2024}
IBBIS. 2024.
\textit{Common Mechanism for DNA Synthesis Screening (commec).}
\url{https://ibbis.bio/commec}

\bibitem{crescendo2024}
Russinovich, M., Salem, A., \& Eldan, R. 2024.
\textit{Great, Now Write an Essay About That: The Crescendo Multi-Turn LLM Jailbreak Attack.}
arXiv:2404.01833.

\end{thebibliography}

\end{document}
